\section{英語3(川村)}
\noindent 【1】AとBの英文について，英文の意味が通るよう，空所に入れるのに適当なものを(あ)$\sim$(く)からひとつずつ選び，記号で答えなさい．
文頭にくるものも小文字で示しています．(1点$\times$16)\\
A:\\
The news of Dr. Nakamura's death in 2019 was a great ( 1 ) to many people, both in ( 2 ) and Japan. 
However, it was ( 3 ) hard for ( 4 ) students studying in Japan, since many of them had been ( 5 ) to study in Japan by Dr. Nakamura himself. 
These students are ( 6 ) to continue Dr. Nakamura's work and to spread his way of thinking and his ( 7 ) approach to getting things done. 
They understand that planning is important, but that doing is ( 8 ) more important. 

\bigskip
  \begin{tabular}{llll}
  (あ) Afghan & (い) Afghanistan & (う) determined & (え) even\\
  (お) inspired & (か) particularly & (き) practical & (く) shock\\
  \end{tabular}
\bigskip\\
B:\\
The problem of food waste will not be ( 1 ) overnight. Laws, ( 2 ) and habits all need to change. 
( 3 ) that food waste is a big ( 4 ) to global warming is already rising, however.
The ( 5 ) of new eco-tech companies may just help to push the change a little faster by making ( 6 ) that ( 7 ) ( 8 ) food that is left over is thrown away. 

\bigskip
  \begin{tabular}{llll}
  (あ) all & (い) appearance & (う) attitudes & (え) awareness\\
  (お) contributor & (か) not & (き) solved & (く) sure\\
  \end{tabular}
\bigskip\bigskip

\noindent 【2】空所に入る語を語群から一つずつ選びなさい．文頭にくるものも小文字で示しています．(1点$\times$20)
\begin{enumerate}
  \item Sending a (\quad\quad) on the network is very easy. 
  \item The system works with only a (\quad\quad) of batteries.
  \item The goal of the race is in the (\quad\quad) of the city. 
  \item The program stopped because of an (\quad\quad).
  \item What was the (\quad\quad) number of engineers who took part in the project? 
  \item There is a red (\quad\quad) in the Japanese flag. 
  \item Put the tools in the right (\quad\quad). 
  \item The (\quad\quad) is too small to show all the information. 
  \item Put the (\quad\quad) through the tube. 
  \item Water changes into (\quad\quad) when it's heated. 
  \item You need a (\quad\quad) when you fly over the ocean. 
  \item The company's (\quad\quad) is to sell more than ten thousand cars a month. 
  \item The (\quad\quad) is driven by the motor. 
  \item If water becomes (\quad\quad), it's called ice. 
  \item Such big waves cannot be seen under (\quad\quad) conditions. 
  \item (\quad\quad) one is one less than zero. 
  \item The pictures are shown in (\quad\quad) of date. 
  \item This program controls the (\quad\quad) of the robot. 
  \item Many people think English is an (\quad\quad) language. 
  \item A hundred thousand (\quad\quad) of the book sold in a month. 
\end{enumerate}
\medskip
\begin{tabular}{|l l l l l|}
  \hline
  center&circle&compass&copies&error\\
  gas&international&message&minus&motion\\
  normal&order&pair&position&pump\\
  screen&solid&target&total&wire\\ 
  \hline
\end{tabular}
\bigskip\bigskip

\noindent 【3】次の表は英単語とその意味についてまとめたものである．(1)$\sim$(5)の空所に入る英語1語を答えなさい．
与えられているアルファベットから始める単語を答えること．(1点$\times$5)\bigskip\\
\begin{tabular}{|l|l|}
  \hline
  \multicolumn{1}{|c|}{英単語}&\multicolumn{1}{|c|}{意味} \\\hline
  consumer&a (1) (p\qquad) who buys goods or uses services \\\hline
  drought&a long period of time with no (2) (\qquad) \\\hline
  (3) (i\qquad)&a system to supply water to farms through pipes or canals \\\hline
  (4) (m\qquad)&things that are used for making something \\\hline
  well&a deep hole in the ground from which people take (5) (w\qquad) \\
  \hline
\end{tabular}
\bigskip\bigskip

\noindent 【4】日本語に合うように空所に入る英語1語をそれぞれ答えなさい．与えられているアルファベットから始まる単語を答えること．(1点$\times$10)
\begin{enumerate}
  \item 彼らはその仕事を完成させるのに1か月かかった．
  \par They took one month to (c\qquad) the work. 
  \item レストランの入口には，「上着とネクタイをお召しください」と書いてあった．
  \par By the entrance of the restaurant, it said, "A jacket and tie are (r\qquad)." 
  \item 台風だったので私たちは一日中テレビを見ながら家にいた．
  \par We stayed home, (w\qquad) TV all day because of the typhoon. 
  \item 私は彼にここに留まるように説得を試みたが，故郷に向けて出て行ってしまった．
  \par I tried to persuade him (t\qquad) stay here, but he left for his hometown. 
  \item 質問してもよろしいでしょうか．
  \par Would you mind (m\qquad) asking a question? 
  \item 私は家族にいまだに子どもだと見なされていることが嫌だ．
  \par I hate to be still (c\qquad) a child by my family.  
  \item 誰だって自分の国が戦争に巻き込まれるのを好まない．
  \par No one wants their own country to be (i\qquad) in war. 
  \item この国ではドラッグは禁止されている．
  \par Drugs are (b\qquad) in this country. 
  \item ビルが今日学校を休んだはずがない．
  \par Bill (c\qquad) have been absent from school today. 
  \item 私は日曜日いつも暇というわけではない．
  \par I am (n\qquad) always free on Sundays. 
\end{enumerate}
\bigskip\bigskip

\noindent 【5】日本語に合うよう【\quad\quad】内の語句を並び替えたとき，2番目と5番目にくるものをそれぞれ記号で答えなさい．(2点$\times$5：完答)
\begin{enumerate}
  \item その店は多くのお客で混んでいた．
  \par [ (あ) a lot\quad(い) customers\quad(う) with\quad(え) of\quad(お) crowded], the shop was doing well.
  \item 毎朝犬を散歩に連れて行くのはあなたの義務よ，ベン．
  \par Ben, [ (あ) walk\quad(い) it's\quad(う) our\quad(え) duty\quad(お) to\quad(か) your\quad(き) dog] every morning. 
  \item 彼女は夫がスポーツカーを買ってくれるのを夢見ている．
  \par She is [ (あ) a sports car\quad(い) her husband\quad(う) of\quad(え) her\quad(お) dreaming\quad(か) buying]. 
  \item この状況はあと3年は続くと見積もられている．
  \par [ (あ) last\quad(い) that\quad(う) is\quad(え) estimated\quad(お) this situation\quad(か) will\quad(き) it] another three years. 
  \item 彼女はずっと，目を閉じたままその問題について考えていた．
  \par She was thinking about the problem [ (あ) with\quad(い) for\quad(う) her\quad(え) closed\quad(お) eyes] a long time. 
\end{enumerate}
\bigskip\bigskip

\noindent 【6】各文の英文の内容が一致するよう空所に英語一語をそれぞれ入れなさい．(1点$\times$4)
\begin{enumerate}
  \item She told me to use a hairdryer to tear off the sticker. 
  \par She told me to (\qquad)(\qquad)(\qquad) a hairdryer to tear off the sticker. 
  \item We showed our thanks for her advice.
  \par We (\qquad)(\qquad)(\qquad) her advice. 
  \item These plastic bottles can be made into clothes.
  \par Clothes can (\qquad)(\qquad)(\qquad) these plastic bottles. 
  \item How can I remove this dirt?
  \par How can I (\qquad)(\qquad)(\qquad) this dirt?
\end{enumerate}
\bigskip\bigskip

\noindent 【7】会話の空所に入れるものとして最も適当なものを，(あ)$\sim$(お)から一つずつ選びそれぞれ記号で答えなさい．(1点$\times$5)\medskip\\
\begin{tabular}{l l}
  Nana:& I've never had Afghan food before.\\
  Ali:& There aren't so many Afghan restaurants in Tokyo. ( 1 ) They serve real traditional food from Afghanistan. \\
  Nana:& Living in Japan, you must miss the food.\\
  Ali:& ( 2 ) It's Palau, our national food. It's steamed rice with goat meat, nuts and fruit. \\
  Nana:& Hmm.... it looks really good.\\
  Ali:& Give it a try....\\
  Nana:& Oh, this is good.\\
  Ali:& I'm glad you like it. As you know, Afghanistan is an agricultural country. ( 3 ) \\
  Nana:& Really? ( 4 ) You need a lot of farmland and water systems for farming then? \\
  Ali:& Yes, but dry weather and floods have greatly damaged the farmlands. That's why there is not enough food,\\
  &and people are suffering from hunger. ( 5 )\\ 
  Nana:& That's terrible.\\
\end{tabular}
\bigskip\par
(あ) But this place is great.\par
(い) In fact, nearly half of Afghan children can't get enough food to eat.\par
(う) My image of Afghanistan is a war zone with dry lands.\par
(え) Oh yes, especially this first dish I ordered.\par
(お) Vegetables, rice, fruits and nuts are important products there. \\

\newpage
\noindent 【8】次の会話文を読み，問いに答えなさい．(1点$\times$5)\medskip\\
\begin{tabular}{l l}
  Hanako:& Steve, I like your T-shirt.\\
  Steve:& Thanks. Actually, I have this shirt in several colors. I have a rainbow in my closet.\\
  Hanako:& What do you mean?\\
  Steve:& I have seven different colors of T-shirts of the same design. Thanks to fast fashion shops, clothes are really\\
  &cheap. So, I can afford to buy several of them.\\
  Hanako:& But you don't need so many T-shirts, do you?\\
  Steve:& Well, I want to look fashionable.\\
  Hanako:& I'm sorry, but you're producing a lot of waste. I don't like the idea of buying so many clothes and throwing\\
  &them away too easily. It's bad for the environment.\\
  Steve:& I don't believe my T-shirts alone are doing serious damage to the earth. It's just a few shirts, you know,\\
  &for fashion.\\
  Hanako:& Come on, Steve! Being eco-friendly is much more in fashion now. Your way of thinking is not fashionable\\
  &at all! \\
\end{tabular}
\medskip\par
\begin{enumerate}
  \item What is Steve concerned about?\par
    \begin{tabular}{l l}
      (あ) Finding a rainbow.&(い) Environmental problems.\\
      (う) Being fashionable.&(え) Cleaning his room.\\
    \end{tabular}
  \bigskip
  \item What does Steve mean when he says, "I have a rainbow in my closet"?\par
    \begin{tabular}{l l}
      (あ) He has clothes with a rainbow.&(い) He has many colorful clothes.\\
      (う) He has only seven T-shirts in his closet.&(え) He likes rainbow colors.\\
    \end{tabular}
  \bigskip
  \item Why does Steve say he can afford to buy many T-shirts?\par\:
      (あ) Because his family is rich and he has a lot of money.\par\:
      (い) Because fast fashion shops sell cheap clothes.\par\:
      (う) Because he wants to wear all seven T-shirts at once.\par\:
      (え) Because he often gives his clothes to his friends.
  \bigskip
  \item What does Hanako think about Steve's habit of buying many T-shirts?\par
    \begin{tabular}{l l}
      (あ) It is a good way to look fashionable.&(い) It helps the environment.\\
      (う) It produces a lot of waste.&(え) It is very eco-friendly.\\
    \end{tabular}
  \bigskip
  \item What does Hanako say is fashionable now?\par
    \begin{tabular}{l l}
      (あ) Wearing rainbow T-shirts.&(い) Being eco-friendly.\\
      (う) Shopping at fast fashion stores.&(え) Throwing away old clothes.\\
    \end{tabular}
\end{enumerate}

\newpage
\noindent 【9】英文を読み，問いに答えなさい．(13点)\par\qquad
I was a university student when I first visited the construction site of Dr. Nakamura's canal. To dam up the river, he was using wire mesh bags filled with stones from the surrounding area instead of concrete.
When I asked him why, he explained how it is better to make use of the natural resources around us.\par\qquad
(1) \underbar{Every time I visited the construction site, what inspired me the most was Dr. Nakamura himself} \\\underbar{working as one of the workers.}
He operated the heavy machinery and carried heavy loads on his back.
He wore the same clothes as the local residents and spoke the local language fluently, so it was often difficult to find him among all the workers.
His diligence and sincerity (2)[ and / him / made / people / trust ] follow him. 
He inspired people.\par\qquad
After meeting Dr. Nakamura, I began to dream of studying in Japan, and in 2017, my dream came true.
I studied international agricultural development, (3) \underbar{mainly focusing on how to make harvested fruits or vegetables}\\\underbar{last longer}.\par\qquad
The last time I saw the doctor was in September 2018.
He encouraged me.
He said he was happy with me focusing on agricultural development.
(4) \underbar{After I completed my master's degree}, I came back to Jalalabad to teach agriculture at the university.
I want to promote Dr. Nakamura's way of thinking and help my hometown to develop.
\bigskip
\begin{enumerate}
  \item 川をせき止めるために，中村医師が使っていたものは何ですか．第1段落で述べられている情報を，過不足なく日本語でまとめなさい．(3点)\vspace{1.8cm}
  \item 下線部(1)を日本語に直しなさい．(4点)\vspace{1.8cm}
  \item (2)の[\qquad]内の語を，意味が通るよう並べ替えなさい．(2点)\vspace{1.8cm}
  \item 下線部(3)と(4)をそれぞれ次のように言い換えたとき，空所に入る語を答えなさい．(2点$\times$2：完答)
  \par (3) and mainly (\qquad)(\qquad)(\qquad) to make harvested fruits or vegetables last longer.
  \par (4) (\qquad)(\qquad)(\qquad) degree,
\end{enumerate}

\newpage
\noindent 【10】英文を読み，問いに答えなさい．(12点)\par\qquad
Reducing food waste will require a lot of effort by governments and individuals.
Less food waste is an important goal because it can lead to more efficient land use and better water resource management.
That's good for the Earth.
Also, many people on the planet do not get enough to eat.
According to the United Nations, the number is more than 820 million people.
Reducing food waste is a way to show respect for food, the farmers who produce it, the natural resources that are used to produce it, and the people who don't get enough of it.\par\qquad
Governments around the world are trying to create new laws that reduce food waste.
One of the first to do so was France.
In 2016, the French government passed a law that banned supermarkets from throwing away food.
The law required large supermarkets to work together with food rescue organizations to donate surplus food to people in need.
The government of Japan has also been considering changes to best-before date laws and laws regarding the donation of food items.
However, it's not enough to make a law.
Someone has to bring the food to the people who need or want it.\par\qquad
In recent years, new start-up companies have been trying to come up with ways to use technology to help with the problem of food waste.
In particular, they are trying to find a way to match leftover food with anyone who would like it.
One such company began operating in Denmark.
They produce an app, which is now available in many other countries, especially in Europe.
The idea is simple: it connects app users to supermarkets, restaurants, bakeries, and other stores that have unsold food.
App users get great discounts on food. The businesses get to reduce their waste, but they also get potential new customers to try their food.
Both businesses and customers contribute to a better environment.
Recently, several similar apps have appeared in Japan as well.
They make it simple for consumers in Japan to help to reduce food waste.
\bigskip
\begin{enumerate}
  \item 以下は本文中に述べられている「食品廃棄を減らす意義」について日本語でまとめたものである．それぞれの空所に入る日本語を答えなさい．(1点$\times$3)
  \par 食品廃棄を減らすことは，(\ding{172})につながる．また，食品廃棄を減らすことにより，食料，食料を生産する農家，食料を生産するために使用される天然資源，そして(\ding{173})に対して(\ding{174})．\bigskip
  \item フランスで2016年に制定された法律について，本文中に述べられている情報を，過不足なく日本語でまとめなさい．(3点)\vspace{1.8cm}
  \item デンマークで開発されたアプリについて，本文中に述べられている情報を，過不足なく日本語でまとめなさい．(6点)
\end{enumerate}